\chapter{1985年全国卷（文）}

\section{单选题}

\begin{problem}

如果正方体 \(ABCD - A'B'C'D'\) 的棱长为 \(a\)，那么四面体 \(A' - ABD\) 的体积是（\quad）

\choices
    {\(\frac{a^3}{2}\)}
    {\(\frac{a^3}{3}\)}
    {\(\frac{a^3}{4}\)}
    {\(\frac{a^3}{6}\)}

\begin{answer}
D
\end{answer}

\begin{solution}
底面三角形 \(ABD\) 是正方形 \(ABCD\) 的一半，所以
\(S_{\triangle ABD}=\frac{1}{2}a^2\). 点 \(A'\) 在底面 \(ABCD\) 上点 \(A\) 的正上方，故四面体 \(A'-ABD\) 的高为 \(A'A=a\). 因此
\(V_{A'-ABD}=\frac{1}{3}S_{\triangle ABD}\cdot A'A=\frac{1}{3}\cdot\frac{a^2}{2}\cdot a=\frac{a^3}{6}\).

所以应选 D.
\end{solution}
\end{problem}

\begin{problem}

\(\tan x = 1\) 是 \(x = \frac{5}{4}\pi\) 的（\quad）

\choices
    {必要条件}
    {充分条件}
    {充分必要条件}
    {既不充分又不必要的条件}

\begin{answer}
A
\end{answer}

\begin{solution}
令命题 \(P\) 为“\(\tan x=1\)”，命题 \(Q\) 为“\(x=\frac{5}{4}\pi\)”. 当 \(Q\) 成立时，
\(\tan x=\tan\frac{5}{4}\pi=1\)，所以 \(Q\) 能推出 \(P\)，即 \(P\) 是 \(Q\) 的必要条件.

但 \(P\) 不能推出 \(Q\)，因为 \(x=\frac{1}{4}\pi\) 时 \(\tan x=1\)，而 \(x\ne\frac{5}{4}\pi\). 所以“\(\tan x=1\)”是“\(x=\frac{5}{4}\pi\)”的必要而非充分条件，应选 A.
\end{solution}
\end{problem}

\begin{problem}

设集合 \(X = \{0,1,2,4,5,7\}\)，\(Y = \{1,3,6,8,9\}\)，\(Z = \{3,7,8\}\)，那么集合 \((X \cap Y) \cup Z\) 是（\quad）

\choices
    {\(\{0,1,2,6,8\}\)}
    {\(\{3,7,8\}\)}
    {\(\{1,3,7,8\}\)}
    {\(\{1,3,6,7,8\}\)}

\begin{answer}
C
\end{answer}

\begin{solution}
先求交集：集合 \(X\) 与 \(Y\) 的公共元素只有 \(1\)，故
\(X\cap Y=\{1\}\). 再与 \(Z\) 作并集，得到
\((X\cap Y)\cup Z=\{1\}\cup\{3,7,8\}=\{1,3,7,8\}\).

所以应选 C.
\end{solution}
\end{problem}

\begin{problem}

在下面给出的函数中，哪一个函数既是区间 \(\left(0,\frac{\pi}{2}\right)\) 上的增函数又是以 \(\pi\) 为周期的偶函数（\quad）

\choices
    {\(y=x^{2}\)（\(x\in\R\)）}
    {\(y=|\sin x|\)（\(x\in\R\)）}
    {\(y=\cos 2x\)（\(x\in\R\)）}
    {\(y=\e^{\sin 2x}\)（\(x\in\R\)）}

\begin{answer}
B
\end{answer}

\begin{solution}
逐项检查两个性质. 函数 \(y=x^2\) 虽为偶函数且在 \(\left(0,\frac{\pi}{2}\right)\) 上递增，但 \(f(0)=0\ne\pi^2=f(\pi)\)，不是以 \(\pi\) 为周期的函数.

函数 \(y=|\sin x|\) 满足 \(|\sin(-x)|=|\sin x|\)，所以是偶函数；又有 \(|\sin(x+\pi)|=|\sin x|\)，所以 \(\pi\) 是它的周期. 在 \(\left(0,\frac{\pi}{2}\right)\) 上，\(|\sin x|=\sin x\)，而 \(\sin x\) 严格递增，因此该函数满足题设的两个条件.

函数 \(y=\cos 2x\) 虽为偶函数且以 \(\pi\) 为周期，但在 \(\left(0,\frac{\pi}{2}\right)\) 上其导数为 \(-2\sin 2x<0\)，是减函数. 函数 \(y=\e^{\sin 2x}\) 在 \(x=\frac{\pi}{6}\) 与 \(x=\frac{\pi}{3}\) 处函数值相同，且 \(f(-x)=\e^{-\sin 2x}\) 一般不等于 \(f(x)\)，既不是该区间上的增函数，也不是偶函数.

所以应选 B.
\end{solution}
\end{problem}

\begin{problem}

用 \(1\)、\(2\)、\(3\)、\(4\)、\(5\) 这五个数字，可以组成比 \(20000\) 大，并且百位数不是数字 \(3\) 的没有重复数字的五位数，共有（\quad）

\choices
    {\(96\text{ 个}\)}
    {\(78\text{ 个}\)}
    {\(72\text{ 个}\)}
    {\(64\text{ 个}\)}

\begin{answer}
B
\end{answer}

\begin{solution}
五位数的首位不能为 \(1\)，且用这些数字组成的数首位为 \(2\)、\(3\)、\(4\) 或 \(5\) 时都大于 \(20000\). 先不考虑百位限制，每确定首位后，其余四位有 \(4!=24\) 种排法，共有 \(4\times24=96\) 个.

现在排除百位为 \(3\) 的情形. 首位为 \(2\)、\(4\) 或 \(5\) 时，数字 \(3\) 仍在其余四位中，固定它在百位后，另外三位有 \(3!=6\) 种排法，所以每种首位应排除 \(6\) 个. 首位为 \(3\) 时，数字 \(3\) 已经用作首位，百位不可能为 \(3\)，因此保留全部 \(24\) 个.

所求个数为
\(\left(24-6\right)+24+\left(24-6\right)+\left(24-6\right)=78\).

所以应选 B.
\end{solution}
\end{problem}

\section{解答题}

\begin{problem}

求函数 \(y = \frac{\sqrt{4 - x^2}}{x - 1} \) 的定义域.
\end{problem}

\begin{solution}
根式有意义要求 \(4-x^2\geqslant0\)，即 \(-2\leqslant x\leqslant2\). 同时分母不能为零，故 \(x\ne1\). 所以定义域为 \(\left[-2,1\right)\cup\left(1,2\right]\).
\end{solution}

\begin{problem}

求圆锥曲线 \(3x^{2} - y^{2} + 6x + 2y - 1 = 0\) 的离心率.
\end{problem}

\begin{solution}
配方得 \(3\left(x+1\right)^2-\left(y-1\right)^2=3\)，所以
\(\frac{\left(x+1\right)^2}{1}-\frac{\left(y-1\right)^2}{3}=1\).

这是双曲线的标准方程，其中 \(a^2=1\)，\(b^2=3\). 因此 \(c^2=a^2+b^2=4\)，从而离心率为 \(e=\frac{c}{a}=2\).
\end{solution}

\begin{problem}

求函数 \(y = -x^{2} + 4x - 2\) 在区间 \([0,3]\) 上的最大值和最小值.
\end{problem}

\begin{solution}
将函数化为顶点式：\(y=-\left(x-2\right)^2+2\). 当 \(x=2\) 时，\(y=2\)，且 \(2\in[0,3]\)，所以最大值为 \(2\). 又因为 \(x\in[0,3]\) 时 \(-2\leqslant x-2\leqslant1\)，故 \(0\leqslant\left(x-2\right)^2\leqslant4\)，从而 \(y=2-\left(x-2\right)^2\geqslant-2\). 当 \(x=0\) 时等号成立，所以最小值为 \(-2\).
\end{solution}

\begin{problem}

设 \((3x-1)^6=a_6x^6+a_5x^5+a_4x^4+a_3x^3+a_2x^2+a_1x+a_0\)，求 \(a_6+a_5+a_4+a_3+a_2+a_1+a_0\) 的值.
\end{problem}

\begin{solution}
在恒等式

\[
(3x-1)^6=a_6x^6+a_5x^5+a_4x^4+a_3x^3+a_2x^2+a_1x+a_0
\]

中令 \(x=1\)，左边为 \((3-1)^6=2^6=64\)，右边为 \(a_6+a_5+a_4+a_3+a_2+a_1+a_0\). 所以所求的值为 \(64\).
\end{solution}

\begin{problem}

设 \(\i\) 是虚数单位，求 \((1+\i)^6\) 的值.
\end{problem}

\begin{solution}
因为 \(\i^2=-1\)，所以 \((1+\i)^2=1+2\i+\i^2=2\i\).

于是 \((1+\i)^6=\left((1+\i)^2\right)^3=(2\i)^3=8\i^3=-8\i\).
\end{solution}

\begin{problem}

设 \(S_1=1^2\)，\(S_2=1^2+2^2+1^2\)，\(S_3=1^2+2^2+3^2+2^2+1^2\)，\(\cdots\)，\(S_n=1^2+2^2+3^2+\cdots+n^2+\cdots+3^2+2^2+1^2\)，\(\cdots\). 用数学归纳法证明：公式 \(S_n=\frac{n(2n^2+1)}{3}\) 对所有的正整数 \(n\) 都成立.
\end{problem}

\begin{solution}
令 \(P(n)\) 表示命题 \(S_n=\frac{n(2n^2+1)}{3}\).

当 \(n=1\) 时，\(S_1=1^2=1\)，而 \(\frac{1(2\cdot1^2+1)}{3}=1\)，所以 \(P(1)\) 成立.

假设 \(P(k)\) 成立，即 \(S_k=\frac{k(2k^2+1)}{3}\)，其中 \(k\) 为任意正整数. 根据 \(S_n\) 的定义，\(S_{k+1}\) 比 \(S_k\) 多出 \(k^2\) 和 \((k+1)^2\)，因此

\[
S_{k+1}=S_k+k^2+\left(k+1\right)^2
=\frac{k\left(2k^2+1\right)}{3}+2k^2+2k+1
=\frac{2k^3+6k^2+7k+3}{3}
\]

又因为 \(2k^3+6k^2+7k+3=\left(k+1\right)\left(2\left(k+1\right)^2+1\right)\)，所以
\(S_{k+1}=\frac{\left(k+1\right)\left(2\left(k+1\right)^2+1\right)}{3}\).

所以 \(P(k+1)\) 成立. 由数学归纳法，公式 \(S_n=\frac{n(2n^2+1)}{3}\) 对所有的正整数 \(n\) 都成立.
\end{solution}

\begin{problem}

证明三角恒等式：\(2\sin^4 x+\frac{3}{4}\sin^2 2x+5\cos^4 x-\cos 3x\cos x=2(1+\cos^2 x)\).
\end{problem}

\begin{solution}
由 \(\sin 2x=2\sin x\cos x\) 和 \(\cos 3x=4\cos^3x-3\cos x\)，可得 \(\frac{3}{4}\sin^2 2x=3\sin^2x\cos^2x\) 以及 \(\cos 3x\cos x=4\cos^4x-3\cos^2x\). 因此原式左边为

\[
2\sin^4x+3\sin^2x\cos^2x+\cos^4x+3\cos^2x
\]

又因为 \(\sin^2x+\cos^2x=1\)，所以

\[
2\sin^4x+3\sin^2x\cos^2x+\cos^4x
=\left(\sin^2x+\cos^2x\right)\left(2\sin^2x+\cos^2x\right)
=2\sin^2x+\cos^2x
\]

从而原式左边等于

\[
2\sin^2x+4\cos^2x
=2\left(\sin^2x+\cos^2x\right)+2\cos^2x
=2\left(1+\cos^2x\right)
\]

即得所证恒等式.
\end{solution}

\begin{problem}

解方程：\(\lg (3-x)-\lg (3+x)=\lg (1-x)-\lg (2x+1)\).
\end{problem}

\begin{solution}
由对数的真数必须为正，需同时满足 \(3-x>0\)，\(3+x>0\)，\(1-x>0\)，\(2x+1>0\)，所以定义域为 \(\left(-\frac{1}{2},1\right)\). 在此范围内，原方程可化为
\(\lg\frac{3-x}{3+x}=\lg\frac{1-x}{2x+1}\).

由于常用对数函数在正数范围内是单调函数，故
\(\frac{3-x}{3+x}=\frac{1-x}{2x+1}\)，从而
\(\left(3-x\right)\left(2x+1\right)=\left(3+x\right)\left(1-x\right)\).

展开得 \(-2x^2+5x+3=-x^2-2x+3\)，即 \(x\left(7-x\right)=0\). 所得根为 \(x=0\) 或 \(x=7\)，其中 \(x=7\) 不在定义域内，故原方程的解为 \(x=0\).
\end{solution}

\begin{problem}

解不等式：\(\sqrt{2x+5}>x+1\).
\end{problem}

\begin{solution}
首先由根式有意义得 \(2x+5\geqslant0\)，即 \(x\geqslant-\frac{5}{2}\).

当 \(-\frac{5}{2}\leqslant x<-1\) 时，\(x+1<0\)，而 \(\sqrt{2x+5}\geqslant0\)，所以不等式成立，得到 \(\left[-\frac{5}{2},-1\right)\).

当 \(x\geqslant-1\) 时，两边均为非负数，可以平方，得

\[
2x+5>\left(x+1\right)^2
\quad\Longleftrightarrow\quad
x^2<4
\quad\Longleftrightarrow\quad
-2<x<2
\]

与 \(x\geqslant-1\) 合并，得到 \(\left[-1,2\right)\). 综上，原不等式的解集为 \(\left[-\frac{5}{2},2\right)\).
\end{solution}

\begin{problem}

已知三棱锥 \(V-ABC\) 的三个侧面与底面所成的二面角都是 \(\beta\)，它的高是 \(h\). 求这个棱锥底面的内切圆半径.
\end{problem}

\begin{solution}
设 \(H\) 是顶点 \(V\) 在底面 \(ABC\) 上的射影，过 \(H\) 向底面边 \(AB\) 作垂线，垂足为 \(D\). 则 \(VH\perp AB\)，\(HD\perp AB\)，所以平面 \(VHD\) 垂直于 \(AB\)，且二面角的平面角为 \(\angle VDH=\beta\). 在直角三角形 \(VHD\) 中，\(VH=h\)，故由 \(\tan\beta=\frac{VH}{HD}=\frac{h}{HD}\)，得 \(HD=h\cot\beta\).

对底面另外两边作同样的垂线，可得 \(H\) 到三边的距离都等于 \(h\cot\beta\). 因而 \(H\) 是三角形 \(ABC\) 的内心，底面的内切圆半径为 \(h\cot\beta\).
\end{solution}

\begin{problem}

已知圆 \(C:x^2+y^2+4x-12y+39=0\) 和直线 \(L:3x-4y+5=0\). 求圆 \(C\) 关于直线 \(L\) 的对称的圆的方程.
\end{problem}

\begin{solution}
将圆的方程配方，得

\(\left(x+2\right)^2+\left(y-6\right)^2=1\).

所以圆心为 \(O\left(-2,6\right)\)，半径为 \(1\). 直线 \(L\) 的方程中 \(3^2+\left(-4\right)^2=25\)，且
\(3\left(-2\right)-4\cdot6+5=-25\).

根据点关于直线的对称公式，圆心的对称点 \(O'\left(x',y'\right)\) 满足

\(x'=-2-\frac{2\cdot3\left(-25\right)}{25}=4\).

同理 \(y'=6-\frac{2\cdot\left(-4\right)\left(-25\right)}{25}=-2\).

对称圆的半径不变，故所求圆的方程为 \(\left(x-4\right)^2+\left(y+2\right)^2=1\).
\end{solution}

\begin{problem}

设首项为 \(1\)，公比为 \(q\)（\(q>0\)）的等比数列的前 \(n\) 项之和为 \(S_n\). 又设 \(T_n=\frac{S_n}{S_{n+1}}\)，\(n=1\)，\(2\)，\(\cdots\). 求 \(\displaystyle\lim_{n\to\infty} T_n\).
\end{problem}

\begin{solution}
当 \(0<q<1\) 时，

\(S_n=\frac{1-q^n}{1-q}\)，所以 \(T_n=\frac{S_n}{S_{n+1}}=\frac{1-q^n}{1-q^{n+1}}\longrightarrow1\).

当 \(q=1\) 时，\(S_n=n\)，所以 \(T_n=\frac{n}{n+1}\longrightarrow1\).

当 \(q>1\) 时，

\(T_n=\frac{q^n-1}{q^{n+1}-1}=\frac{1-q^{-n}}{q\left(1-q^{-(n+1)}\right)}\longrightarrow\frac{1}{q}\).

因此 \(\displaystyle\lim_{n\to\infty}T_n=\begin{cases}
1, & 0<q\leqslant1,\\
\dfrac{1}{q}, & q>1
\end{cases}\).
\end{solution}
